Die Rampe ist ein rechtwinkliges Dreieck: $\ell = \lR$ entlang der Rampe, $h = \hR$ hoch, also waagrecht (Pythagoras)
\begin{align}
 b &= \bRF = \bR
\end{align}
Für den Neigungswinkel $\alpha$ gilt damit $\sin\alpha = h/\ell$ und $\cos\alpha = b/\ell$.

\begin{abcliste}
\abc
In Ruhe hält die Haftreibung der Hangabtriebskraft $m\,g\sin\alpha$ das Gleichgewicht, die Normalkraft der Komponente $m\,g\cos\alpha$ in die Rampe hinein:
\begin{align}
 \ssc{F}{R} &= \FRF \\
   &= \m\cdot\g\cdot\frac{\hR}{\lR} \\
   &= \FR \approx \result{\FRP{3}{2}}
\end{align}
\begin{align}
 \ssc{F}{N} &= \FNF \\
   &= \m\cdot\g\cdot\frac{\bR}{\lR} \\
   &= \FN \approx \result{\FNP{3}{3}}
\end{align}

\abc
Das Auto hält, wenn die nötige Reibung höchstens $\mu_\mathrm{H}\,\ssc{F}{N}$ ist:
\begin{align}
 \ssc{F}{R} &\le \mu_\mathrm{H}\,\ssc{F}{N} \\
 \mu_\mathrm{H} &\ge \frac{\ssc{F}{R}}{\ssc{F}{N}} = \muHF \\
   &= \frac{\hR}{\bR} \\
   &= \muH \approx \result{\muHP{0}{2}}
\end{align}
Die Masse kürzt sich weg: Ob ein Auto rutscht, hängt nur von der Steigung und den Reifen ab.
\end{abcliste}
